%------------------------------------------------------------------------------%
\begin{question}
Encontre uma parametrização para o movimento de uma partícula que começa no
ponto \((-2, 0)\) e traça a metade superior do círculo \(x^2+y^2=4\)
\end{question}

%------------------------------------------------------------------------------%
\begin{resolution}{Solução}
  Metade superior do círculo de raio 2 centrado na origem

  Parametrização do círculo unitário centrado na origem
  \begin{align*}
    x      & = \cos(\theta) \\
    y      & = \sin(\theta) \\
    \theta & \in [0, 2\pi)
  \end{align*}
\end{resolution}

%------------------------------------------------------------------------------%
\begin{resolution}{Solução}
  Parametrização do círculo de raio 2 centrado na origem
  \begin{align*}
    x      & = 2\cos(\theta) \\
    y      & = 2\sin(\theta) \\
    \theta & \in [0, 2\pi)
  \end{align*}
\end{resolution}

%------------------------------------------------------------------------------%
\begin{resolution}{Solução}
  Parametrização da metade superior
  \begin{align*}
    x      & = 2\cos(\theta) \\
    y      & = 2\sin(\theta) \\
    \theta & \in [0, \pi]
  \end{align*}
\end{resolution}

%------------------------------------------------------------------------------%
\begin{resolution}{Solução}
  Essa parametrização começa em \((2, 0)\) e termina em \((-2, 0)\)


  Queremos inverter o sentido em que percorremos a trajetória
  \begin{align*}
    x      & = 2\cos(\pi -\theta) =    - 2 \cos(\theta)\\
    y      & = 2\sin(\pi -\theta) = \hpm 2 \sin(\theta)\\
    \theta & \in [0, \pi)
  \end{align*}
\end{resolution}

%------------------------------------------------------------------------------%
